\section{Đẳng biến, không phải bất biến}
\label{sec:ch10-dang-bien}

\index{thiên kiến quy nạp!đẳng biến khác bất biến}%
Hai phát biểu, và cả ngành hay gọi chung bằng một tên.

\tn{Đẳng biến}{equivariance} với phép dịch: dịch đầu vào thì đầu ra dịch theo
đúng như vậy. \tn{Bất biến}{invariance} với phép dịch: dịch đầu vào thì đầu ra
\emph{không đổi}. Zhang viết cả hai ra cạnh nhau ở mục 3.1 bài
mình~\autocite{zhang2019shift}:

\begin{quote}
\enquote{A function F is shift-equivariant if shifting the input equally shifts
the output [...] A representation is shift-invariant if shifting the input
results in an identical representation.}
\end{quote}

Cái thứ nhất là một định lý về phép \tn{tích chập}{convolution}, và mục này đo
nó ra số không tuyệt đối. Cái thứ hai là một phát biểu về cả mạng, và mục này
đo nó ra sai.

Viết bằng ký hiệu thì đẳng biến là một phép hoán vị thứ tự. Gọi \(T_\delta\) là
phép dịch đi \(\delta\) và \(f\) là tầng tích chập:

\begin{equation}
\label{eq:ch10-dang-bien}
f(T_\delta x) = T_\delta f(x).
\end{equation}

Chương~\ref{ch:transformer} đã đo một phát biểu đúng dạng này một lần rồi, cho
phép hoán vị thay vì phép dịch, và cách đo giống hệt: chạy hai vế, trừ nhau,
nhìn số.

\subsection*{Một tầng, đo ra số không}

\index{thiên kiến quy nạp!đo đẳng biến của một tầng}%
\eqref{eq:ch10-dang-bien} là một đẳng thức về số học chứ không phải về trọng
số, nên nó đúng ngay tại khởi tạo và không cần huấn luyện gì. Đó cũng là lý do
con số dưới đây là số không chứ không phải số nhỏ.

\begin{measured}{sai khác hai vế của \eqref{eq:ch10-dang-bien}, dịch một pixel}
\begin{minted}{text}
padding     max |f(shift x) - shift f(x)|
circular    0.000e+00
zeros       1.772e+00
\end{minted}

\begin{minted}{text}
  whole feature map          1.772e+00
  interior, borders dropped  0.000e+00
\end{minted}

Đo bằng \code{experiments/ch10\_equivariance.py} tại tag \code{ch10}.
\end{measured}

Với padding vòng thì hai vế bằng nhau tới từng bit. Với padding không thì
không, và hàng thứ hai nói chỗ hỏng nằm ở đâu: bỏ hai pixel viền ra, phần trong
cũng đúng số không.

Nên cách đọc đúng không phải là padding không làm đẳng biến yếu đi. Nó phá đẳng
biến trên một khung hai pixel và không phá chỗ nào khác. Đó là điều
\tn{tính cục bộ}{locality} tiên đoán: một unit chỉ hỏng khi
\tn{trường tiếp nhận}{receptive field} của nó thò ra ngoài ảnh, và chỉ những
unit sát mép mới thò ra.

\subsection*{Chỗ tính chất ấy mất: phép giảm mẫu}

\index{thiên kiến quy nạp!giảm mẫu phá đẳng biến ở phép dịch lẻ}%
Ràng buộc thứ ba của bài 1998, \tn{giảm mẫu}{subsampling}, là chỗ phải trả giá.

Một tầng giảm mẫu hai lần thì đầu ra chỉ còn nửa độ phân giải, nên dịch đầu vào
đi 2 phải cho đầu ra dịch đi 1. Đem so với đầu ra dịch đi 2 là so sai đại
lượng, và tôi đã so sai đúng như vậy ở bản đầu; mục~\ref{sec:ch10-con-lai-gi}
kể chỗ đó.

\begin{measured}{conv rồi MaxPool2d(2), so đúng tỷ lệ}
\begin{minted}{text}
input shift  output shift  error
0            0             0.000e+00
2            1             0.000e+00
4            2             0.000e+00
6            3             0.000e+00
8            4             0.000e+00
\end{minted}

Đo bằng \code{experiments/ch10\_equivariance.py} tại tag \code{ch10}.
\end{measured}

Đúng số không ở mọi phép dịch chẵn. Còn phép dịch lẻ thì không có đầu ra nào để
đem so cả, nên bảng dưới lấy giá trị tốt nhất trên mọi phép dịch đầu ra.

\begin{measured}{phép dịch lẻ: chỗ tốt nhất cũng không gần}
\begin{minted}{text}
input shift  best output shift  residual  residual/scale
1            -2                 2.327e+00 0.9439
3            -3                 2.258e+00 0.9159
5            -2                 2.258e+00 0.9159
\end{minted}

Biên độ đầu ra là 2.4651.

Đo bằng \code{experiments/ch10\_equivariance.py} tại tag \code{ch10}.
\end{measured}

Cột cuối là chỗ đọc ra. Chỗ khớp nhất mà một phép dịch lẻ đạt được vẫn lệch hơn
chín phần mười biên độ của chính đầu ra. \textbf{Đây không phải một sự suy giảm
tăng dần.} Tầng ấy đúng tuyệt đối ở một nửa số phép dịch và không nói được gì
về nửa còn lại, và hình~\ref{fig:ch10-chan-le} là lý do.

% The picture measures 268.01pt against this book's 441.02pt measure, so it has
% 173pt to spare and could carry two more sampling rows if a later revision
% wants them. Measured with \savebox per quyet dinh 67; see the comment at the
% chapter's first figure for why the built page and the gate are both blind
% here, and for the savebox-naming trap.
\begin{figure}[htbp]
  \centering
  % Vi sao "giu moi mau thu hai" (stride 2) dang bien voi dich chan nhung khong
% dang bien voi dich le: day la ket qua do duoc cua chuong (bias quy nap cua
% CNN). Dung "dang bien" chu khong dung "tuong duong", vi "tuong duong" da
% mang nghia equivalent trong sach nay - xem quyet dinh 70.
% Mot day 10 vi tri dau vao (0..9), moi vi tri la mot cham; mau GIU (tron
% dam, cac vi tri chan) va mau BO (tron vien nhat, cac vi tri le) la CO DINH
% - giong het nhau o ca ba hang - chi mot DAC TRUNG duoc danh dau (hinh thoi)
% la doi cho giua cac hang, dung vi tri no doi toi moi la dieu thay doi.
% Hang 1 (dich 0, goc): dac trung o vi tri 4 (chan) -> roi vao mau giu, la
% dau ra y_2.
% Hang 2 (dich +2, chan): dac trung sang vi tri 6, van roi vao mau giu, dau
% ra y_3 = y_2 dich dung 1 buoc - tuong duong dich hoan hao.
% Hang 3 (dich +1, le): dac trung sang vi tri 5, roi dung vao mau BO - khong
% con o dau ra nao chua dac trung nay nua. Day khong phai "sai lech mot
% chut" ma la mot tap mau hoan toan khac (sublattice khac).
% Khong ve: dau vao 2D, nhieu kenh, gia tri thuc cua dac trung - hinh chi
% can VI TRI cua no.
% Mono-safe: giu/bo phan biet bang to dam vs vien nhat (khong mau); hang
% khop dung mui ten net lien, hang lech dung net cham + dau gach cheo.
% Do rong DA DO bang \savebox tai cho goi hinh: 268.01pt = 9.42cm, tren mot
% measure 441.02pt. Do la con so co gia tri; dung uoc luong lai bang mat.
\begin{tikzpicture}[
  >= Stealth,
  giu/.style      = {draw, circle, fill = black, minimum size = 3.0mm,
                      inner sep = 0pt},
  bo/.style       = {draw = black!35, circle, fill = white,
                      minimum size = 3.0mm, inner sep = 0pt},
  % Hình thoi tự chế bằng hình vuông xoay 45 độ, vì thư viện shapes.geometric
  % (có sẵn shape "diamond") không nằm trong danh sách thư viện của sách.
  dactrung/.style = {draw = black, rectangle, line width = 1pt,
                      minimum size = 3.8mm, inner sep = 0pt, fill = none,
                      rotate = 45},
  khop/.style     = {->, black!70, line width = 0.8pt},
  lech/.style     = {black!55, densely dotted, line width = 0.8pt},
  ten/.style      = {font = \bfseries\scriptsize},
  nho/.style      = {font = \scriptsize},
  be/.style       = {font = \tiny, black!55},
  ghichu/.style   = {font = \footnotesize, black!60},
]

  % ================= chú giải ba ký hiệu, phía trên cùng =================
  \node[giu] (leg-giu) at (1.1, 5.9) {};
  \node[nho, anchor = west] at (1.25, 5.9) {giữ};
  \node[bo]  (leg-bo)  at (2.6, 5.9) {};
  \node[nho, anchor = west] at (2.75, 5.9) {bỏ};
  \node[dactrung] (leg-dt) at (4.1, 5.9) {};
  \node[nho, anchor = west] at (4.3, 5.9) {đặc trưng};

  % ================= thước vị trí đầu vào, 0..9, cố định cho cả ba hàng ====
  \foreach \k in {0, 1, 2, 3, 4, 5, 6, 7, 8, 9} {
    \node[be] at (\k*0.7, 5.3) {\k};
  }

  % ================= HÀNG 1: dịch 0 (gốc), y = 4.6 =================
  \node[ten, anchor = east] at (-0.3, 4.6) {Dịch 0 (gốc)};
  \foreach \k in {0, 2, 4, 6, 8} {
    \node[giu] (r1-\k) at (\k*0.7, 4.6) {};
  }
  \foreach \k in {1, 3, 5, 7, 9} {
    \node[bo] (r1-\k) at (\k*0.7, 4.6) {};
  }
  \node[dactrung] at (2.8, 4.6) {}; % vị trí 4, chẵn -> rơi vào mẫu giữ

  \draw[khop] (2.8, 4.35) -- (2.8, 4.05);
  \node[nho, anchor = north] at (2.8, 4.0) {\(y_{2}\)};

  % ================= HÀNG 2: dịch +2, chẵn, y = 2.6 =================
  \node[ten, anchor = east] at (-0.3, 2.6) {Dịch \(+2\) (chẵn)};
  \foreach \k in {0, 2, 4, 6, 8} {
    \node[giu] (r2-\k) at (\k*0.7, 2.6) {};
  }
  \foreach \k in {1, 3, 5, 7, 9} {
    \node[bo] (r2-\k) at (\k*0.7, 2.6) {};
  }
  \node[dactrung] at (4.2, 2.6) {}; % vị trí 6, chẵn -> vẫn rơi vào mẫu giữ

  \draw[khop] (4.2, 2.35) -- (4.2, 2.05);
  \node[nho, anchor = north] at (4.2, 2.0) {\(y_{3}\) (\(=y_{2}\) dịch 1 bước)};

  % ================= HÀNG 3: dịch +1, lẻ, y = 0.6 =================
  \node[ten, anchor = east] at (-0.3, 0.6) {Dịch \(+1\) (lẻ)};
  \foreach \k in {0, 2, 4, 6, 8} {
    \node[giu] (r3-\k) at (\k*0.7, 0.6) {};
  }
  \foreach \k in {1, 3, 5, 7, 9} {
    \node[bo] (r3-\k) at (\k*0.7, 0.6) {};
  }
  \node[dactrung] at (3.5, 0.6) {}; % vị trí 5, lẻ -> rơi đúng vào mẫu bị bỏ

  \draw[lech] (3.5, 0.35) -- (3.5, 0.05);
  \draw[black!55, line width = 0.8pt]
    (3.42, -0.03) -- (3.58, 0.13);
  \draw[black!55, line width = 0.8pt]
    (3.42, 0.13) -- (3.58, -0.03);
  \node[nho, anchor = north] at (3.5, -0.05) {không rơi vào ô đầu ra nào};

  % ================= chú giải tóm tắt, dưới cùng =================
  \node[ghichu, anchor = north, align = left, text width = 8.2cm]
    at (2.5, -1.1)
    {Mẫu giữ/bỏ cố định qua mọi phép dịch - chỉ đặc trưng di chuyển. Dịch
     chẵn luôn hạ đúng vào mẫu giữ, nên đầu ra chỉ là đầu ra gốc dịch đi một
     số nguyên bước: đẳng biến đúng tuyệt đối. Dịch lẻ hạ đúng vào mẫu bị
     bỏ: đặc trưng biến mất khỏi đầu ra, không bước dịch nào của đầu ra tái
     tạo lại được nó.};

\end{tikzpicture}

  \caption{Vì sao phép dịch chẵn và phép dịch lẻ khác nhau về bản chất chứ
  không về mức độ. Phép giảm mẫu giữ lại một mẫu trên hai, tức nó chọn một
  trong hai lưới con. Dịch đi 2 thì tập mẫu giữ lại rơi đúng vào cùng lưới con
  ấy và đầu ra chỉ trượt đi một ô. Dịch đi 1 thì nó rơi vào lưới con \emph{kia},
  và không có phép trượt nào trên đầu ra làm hai bên gặp nhau.}
  \label{fig:ch10-chan-le}
\end{figure}

\subsection*{Không phải do max, mà do bỏ mẫu}

\index{thiên kiến quy nạp!chẵn lẻ là tính chất của phép bỏ mẫu}%
Bài của Zhang nói về aliasing, và nó hay bị đọc thành \enquote{max pooling gây
ra chuyện này}. Đem đúng phép đo trên cho ba cách giảm một nửa lưới thì thấy
không phải.

\begin{measured}{ba cách giảm mẫu, cùng một phép đo}
\begin{minted}{text}
layer           even shift  odd shift  scale   odd/scale
MaxPool2d(2)    0.000e+00   2.215e+00  2.429   0.912
AvgPool2d(2)    0.000e+00   1.101e+00  1.343   0.820
Conv stride 2   0.000e+00   2.968e+00  2.427   1.223
\end{minted}

Đo bằng \code{experiments/ch10\_equivariance.py} tại tag \code{ch10}.
\end{measured}

Cả ba đúng tuyệt đối ở phép dịch chẵn, cả ba sai cỡ biên độ tín hiệu ở phép
dịch lẻ. Hàng thứ ba không có phép max nào trong nó. Nên phân đôi chẵn lẻ là
tính chất của việc vứt một mẫu trên hai, chứ không phải của phi tuyến đứng cạnh
nó.

Chồng thêm tầng thì chu kỳ nhân lên.

\begin{measured}{chu kỳ đẳng biến theo số tầng giảm mẫu}
\begin{minted}{text}
pool layers  total stride  exact at shifts of
1            2             2
2            4             4
3            8             8
\end{minted}

Đo bằng \code{experiments/ch10\_equivariance.py} tại tag \code{ch10}.
\end{measured}

Mạng ba tầng pooling trên ảnh 32 pixel chỉ đẳng biến đúng ở bội của 8, tức 4
trong 32 phép dịch ngang.

\textbf{Và chuyện này bài 1989 đã nói ra.} Cùng mục 3.3 vừa trích ở
mục~\ref{sec:ch10-hai-rang-buoc}, về chính phép giảm mẫu hai pixel của
nó~\autocite{lecun1989zipcode}:

\begin{quote}
\enquote{For units in layer H1 that are one unit apart, their receptive fields
(in the input layer) are two pixels apart. Thus the input image is undersampled
and some position information is eliminated.}
\end{quote}

\enquote{Some position information is eliminated} là câu đúng, viết năm 1989.
Cái mất ba mươi năm là đo xem \emph{bao nhiêu}, và câu trả lời là một nửa số
phép dịch, chứ không phải một ít ở mỗi phép dịch.

\subsection*{Vì sao chỗ này không phải chuyện chữ nghĩa}

\index{thiên kiến quy nạp!vì sao phân biệt hai khái niệm}%
Người ta thường nói mạng tích chập \enquote{bất biến với phép dịch}, và ý là
mạng nhận ra con mèo dù nó ở góc nào. Phát biểu ấy sai theo hai đường khác
nhau, và cả hai đều đo được.

Đường thứ nhất là mục vừa rồi: các đặc trưng chỉ đẳng biến ở bội của tổng
stride.

Đường thứ hai sâu hơn và mục~\ref{sec:ch10-anh-that} đo nó trên mô hình đã
huấn luyện thật. Kể cả khi \eqref{eq:ch10-dang-bien} đúng tuyệt đối, cái nó nói
là \emph{bản đồ đặc trưng dịch theo}. Bộ phân loại đứng sau đọc bản đồ ấy đã
phẳng hóa, nên một bản đồ dịch đi một ô là một vector đầu vào khác hẳn với nó.
\tn{Bản đồ đặc trưng}{feature map} đẳng biến không cho ra đầu ra bất biến; nó
chỉ cho ra đầu ra mà một lớp đọc-toàn-cục \emph{có thể} làm cho bất biến nếu
người ta bắt nó làm vậy.

Đó là lý do kiến trúc thật hay đặt một phép gộp toàn cục trước bộ phân loại.
Bất biến không rơi ra từ tích chập; nó phải được xây thêm.
